\section{Related Work}
\label{sec:detailed_related_work}



\paragraph{Video Generation Models.}
Recent advancements in video generation models have emerged from two primary architectures: diffusion-based models \cite{gen2, svd, liu2024sora, blattmann2023align, wang2023lavie, wang2023modelscope, chen2024videocrafter2, khachatryan2023text2video} and autoregressive modeling-based approaches \cite{magvit2, kondratyuk2023videopoet, hong2022cogvideo, villegas2022phenaki}. Among these, diffusion models have garnered significant attention. The model known as SVD \cite{svd}, built on a Latent Diffusion Model (LDM) \cite{ldm}, proposes a three-stage training process for video LDMs: text-to-image pretraining, video pretraining, and video finetuning. Sora \cite{liu2024sora} represents a state-of-the-art in video generation, utilizing a diffusion-transformer architecture with unified training recipes and enhancements in language description processing for video generation. ModelScope \cite{wang2023modelscope} is also a diffusion-based text-to-video model which combines a VQGAN \cite{esser2021taming}, a text-encoder, and a denoising UNet. Another diffusion model, VideoCrafter2 \cite{chen2024videocrafter2}, leverages low-quality videos and high-quality videos to generate high-quality videos. LaVIE \cite{wang2023lavie} is composed of a base text-to-video model, a temporal interpolation model, and a video super-resolution model, indicating that joint image and video training and temporal self-attention with rotary positional embeddings are key components to boost performance. Given the rapid development of video generation technology, an effective evaluation method for the generated videos becomes crucial. Our paper focuses on evaluating text-to-video generation models for their physical commonsense capabilities.


\paragraph{Evaluating Video Generation Models.}
To evaluate the quality of a T2V generative model, Fréchet video distance (FVD) is traditionally used to measure the similarity between real and generated video distributions \cite{unterthiner2018towards,bugliarello2024storybench}. However, FVD has several limitations for assessing physical commonsense including the requirement for a reference video that is difficult to obtain for novel scenes, bias towards video quality, and failure to detect unrealistic motions \cite{brooks2022generating,skorokhodov2022stylegan}. Similarly, CLIPScore \cite{radford2021learning} measures \textit{semantic} similarity between generated video frames and the conditioning text in a shared representation space, making it unsuitable for evaluating physical commonsense in generated videos.

However, there is a growing consensus on the need for more comprehensive metrics to assess the performance of video generation models \cite{huang2023vbench, liu2024evalcrafter, subjectbench, lin2024evaluating}. V-Bench \cite{huang2023vbench} offers a detailed benchmark suite that introduces a hierarchical evaluation protocol, breaking down `video generation quality’ into various granular perspectives. Another framework, EvalCrafter \cite{liu2024evalcrafter}, proposes 17 objective metrics. Despite these advancements, existing methods largely overlook the fundamental aspect of physical commonsense. Unlike static images, videos incorporate a temporal dimension, embedding physical commonsense information across frames. Our research dives into the measurement of physical commonsense \cite{bisk2020piqa} in videos. While prior research \cite{ku2023viescore} presents training-free image-text alignment evaluation using existing large multimodal models, we find that they are insufficient for accurate judgments on \name{}. Hence, we introduce a \model{} auto-evaluator that is trained on diverse generated videos and human annotations for reliable judgments. 

\paragraph{Physics Modeling.}
Simulating physical behaviors of solids and fluids has always been an important and popular topic in computer graphics. For solid materials, the simplest physical model is the long-established rigid body simulation \cite{baraff1997introduction}, where solids are assumed not to deform. Simulation of deformable solids \cite{sifakis2012fem}, on the other hand, takes into account the strain and stress during deformation. To capture more complicated materials, researchers have been proposing increasingly intricate models for different materials, such as metal \cite{o2002graphical}, sand \cite{klar2016drucker}, and snow \cite{stomakhin2013material}. In contrast, most of the common fluids \cite{bridson2015fluid} in daily life can be broadly categorized as inviscid \cite{koschier2020smoothed}, e.g., water and air, and viscous fluids \cite{takahashi2014fast, larionov2017variational}, e.g., honey and oil. Additionally, an orthogonal research direction is to accurately, efficiently, and robustly model contact and interaction between different materials. These include solid-solid \cite{li2020incremental, li2020codimensional}, solid-fluid \cite{batty2007fast, xie2023contact}, and fluid-fluid interactions \cite{muller2005particle}. Further, recent advancements in computer vision have started exploring incorporating physics priors into various 3D-aware generation tasks to enhance physical plausibility, such as human animation \cite{yuan2023physdiff, shimada2020physcap, xu2023interdiff} and 3D/4D generation \cite{mezghanni2021physically, xie2023physgaussian, zhang2024physdreamer}. \tx{However, these approaches often depend on high-quality 3D reconstructions from multi-view images. Some efforts \cite{liuphysgen} have also integrated physics-based simulations into video generative models, but the simulations are performed in 2D space due to the lack of 3D information, resulting in limited dynamics.} In this work, instead of generating, we focus on identifying whether the generated video adheres to physical laws.


